\providecommand{\Res}[2]{\qty[round-mode=figures,round-precision=3]{\fpeval{#1}}{#2}}
\providecommand{\Val}[1]{\num[round-mode=figures,round-precision=4]{\fpeval{#1}}}
\providecommand{\ResSci}[2]{\qty[round-mode=figures,round-precision=3,exponent-mode=scientific]{\fpeval{#1}}{#2}}
\subsubsection*{Question a) Avant le croisement}
\begin{enumerate}[label=\arabic*)]
    \item \textbf{Variables connues et recherchées :} $f = \qty{\v{f}}{Hz}$, $v_s = \qty{\v{vs}}{m/s}$ et $v_o = \qty{\v{vo}}{m/s}$ (l'un vers l'autre), $v = \qty{\v{v}}{m/s}$ ; variable recherchée : $f'$
    \item \textbf{Équations pertinentes :}
    \[ f' = f\,\frac{v \pm v_o}{v \mp v_s} \]
    \item \textbf{Isolation de la variable recherchée :} les deux s'approchent : $+$ au numérateur, $-$ au dénominateur.
    \[ f' = f\,\frac{v + v_o}{v - v_s} \]
    \item \textbf{Résolution numérique :}
    \[ f' = \num{\v{f}} \cdot \frac{\num{\v{v}} + \num{\v{vo}}}{\num{\v{v}} - \num{\v{vs}}} = \Res{\v{f}*(\v{v}+\v{vo})/(\v{v}-\v{vs})}{Hz} \]
    \item \textbf{Réponse :} elle perçoit \textbf{\Res{\v{f}*(\v{v}+\v{vo})/(\v{v}-\v{vs})}{Hz}}.
\end{enumerate}

\subsubsection*{Question b) Après le croisement}
\begin{enumerate}[label=\arabic*)]
    \item \textbf{Variables connues et recherchées :} mêmes données, les deux s'éloignent ; variable recherchée : $f''$
    \item \textbf{Équations pertinentes :}
    \[ f'' = f\,\frac{v - v_o}{v + v_s} \]
    \item \textbf{Isolation de la variable recherchée :} l'équation est déjà isolée.
    \item \textbf{Résolution numérique :}
    \[ f'' = \num{\v{f}} \cdot \frac{\num{\v{v}} - \num{\v{vo}}}{\num{\v{v}} + \num{\v{vs}}} = \Res{\v{f}*(\v{v}-\v{vo})/(\v{v}+\v{vs})}{Hz} \]
    \item \textbf{Réponse :} elle perçoit \textbf{\Res{\v{f}*(\v{v}-\v{vo})/(\v{v}+\v{vs})}{Hz}}.
\end{enumerate}

\subsubsection*{Question c) Même vitesse, même direction}
L'auto suit le camion : la source s'éloigne de l'observateur ($v + v_s$ au dénominateur) et l'observateur s'approche de la source ($v + v_o$ au numérateur). Avec $v_o = v_s$, les deux facteurs se compensent : elle percevrait \textbf{\qty{\v{f}}{Hz}}, la fréquence émise (aucune vitesse relative, air immobile).
